\documentclass[aps]{revtex4-2}%

\usepackage{algorithm}
\usepackage{algpseudocode}
\begin{document}
\begin{algorithm}
\caption{An algorithm with caption}\label{alg:cap}
\begin{algorithmic}
\Require $n \geq 0$
\Ensure $y = x^n$
\State $y \gets 1$
\State $X \gets x$
\State $N \gets n$
\While{$N \neq 0$}
\If{$N$ is even}
    \State $X \gets X \times X$
    \State $N \gets \frac{N}{2}$  \Comment{This is a comment}
\ElsIf{$N$ is odd}
    \State $y \gets y \times X$
    \State $N \gets N - 1$
\EndIf
\EndWhile
\end{algorithmic}
\end{algorithm}

Due to theorem~\ref{theo:Noisy_reduction},  it is sufficient to focus on graph states under single qubit incoherent noise. It is well known that graph state is generated by $CZ$ gates operate on $\ket{+}^{\otimes n}$ state
\begin{align}
\ket{G} = \prod_{(i,j) \in E} CZ_{ij} \ket{+}^{\otimes |V|},    
\end{align}
The first step of 'back-propagation' of graph states is decomposition of incoherent noise into components that commute with $CZ$ gates. This is proved by the theorem below
\begin{lemma}\label{lemma:commute_channel}
    Incoherent channel that commute with $CZ$ gates must has the form   \(\Lambda_{\Gamma_0}\left(\frac{1}{2}[\mathbb{I} + \mathbf{w} \cdot  \vec\sigma]\right) = \frac{1}{2}[\mathbb{I} + (\mathbf{t} + \Gamma\mathbf{w}) \cdot  \vec\sigma]\) with $\mathbf{t} = (0,0,0), \quad  \Gamma = \text{diag}(\Gamma_0,\Gamma_0,1)$. 
\end{lemma}
\begin{proof}
    It is sufficient to show that $\mathbb{I}\otimes \Lambda(CZ\cdot \rho \cdot CZ) = CZ \cdot  \mathbb{I}\otimes\Lambda(\rho) \cdot CZ$, where $\rho$ is two qubit state, thus can be decomposed as a linear combination of $\{\sigma_0 = \mathbb{I}, \sigma_1,\sigma_2,\sigma_3\}^{\otimes 2}$   
    $CZ \cdot (\sigma_i \otimes \sigma_j) \cdot CZ = (\sigma_i \sigma_3^{\tau(j)}) \otimes (\sigma_3^{\tau(i)}\sigma_j)$ , where $\tau(1)=\tau(2) = 1, \tau(0)=\tau(3) = 0$. 
    Then we have
    \begin{align}
        \mathbb{I}\otimes \Lambda\left(CZ\cdot (\sigma_i \otimes \sigma_j) \cdot CZ\right)  &= (\sigma_i \sigma_3^{\tau(j)}) \otimes \Lambda(\sigma_3^{\tau(i)}\sigma_j),\\
        &= 2^{-1}\sum_{k=0}^3 (\sigma_i \sigma_3^{\tau(j)}) \otimes \sigma_k \cdot \text{tr}\left[ \Lambda(\sigma_3^{\tau(i)}\sigma_j) \sigma_k\right],\\
        CZ \cdot  \mathbb{I}\otimes\Lambda(\sigma_i \otimes \sigma_j) \cdot CZ&= 2^{-1}\sum_{ k =0}^3 ( \sigma_i\sigma_3^{\tau(k)}) \otimes (\sigma_3^{\tau(i)} \sigma_k) \cdot \text{tr}\left[\Lambda(\sigma_j)\sigma_k\right] \\
        &= 2^{-1}\sum_{k'= 0}^3(\sigma_i \sigma_3^{\tau ( u(i,k'))}) \otimes \sigma_{k'} \cdot \text{tr}\left[\sigma_3^{\tau(i)} \Lambda(\sigma_j) \sigma_{k'}\right],
    \end{align}
    where $u: (i,k') \to k$ is a map $\{0,1,2,3\}^2 \to \{0,1,2,3\}$,  such that  $  \sigma_{u(i,k')}\propto\sigma_{3}^{\tau(i)}\sigma_{k'}  $.   
    Then it can be concluded that for all $i,j,k \in \{0,1,2,3\}$, 
     \begin{align}
        \text{tr}\left[ \Lambda(\sigma_3^{\tau(i)}\sigma_j) \sigma_k\right] = \text{tr}\left[\sigma_3^{\tau(i)} \Lambda(\sigma_j) \sigma_{k}\right].
    \end{align}
    This implys that $\Lambda(\sigma_3\sigma_j ) = \sigma_3\Lambda(\sigma_j)$ must be satisfied for $j \in \{0,1,2,3\}$.
    Besides, the identity $\tau(j) = \tau(u(i,k))$ should be satisfied unless 
    \begin{align}
        \text{tr}\left[ \Lambda(\sigma_3^{\tau(i)}\sigma_j) \sigma_k\right] = \text{tr}\left[\sigma_3^{\tau(i)} \Lambda(\sigma_j) \sigma_{k}\right] = 0.
    \end{align}
    For inputs $k \in \{0,3\}$, it must have $u(i,k) \in \{0,3\}$ either because different value of $i$ only change the $\sigma_3$ component in $\sigma_{u(i,k)}$  . 
    Thus the identity $\tau(j) = \tau(u(i,k))$ can be satisfied if and only if $j \in \{0,3\}$. 
    This means that the value $\text{tr}\left[\sigma_3^{\tau(i)} \Lambda(\sigma_j) \sigma_{k}\right]$ should always vanish at  $j \in \{1,2\}$.
    Then the condition above is equivalent to $\Lambda(\sigma_1), \Lambda(\sigma_2)$ has no $\sigma_0, \sigma_3$ components. 
    Similarly, for $k \in \{1,2\}$, we conclude $\Lambda(\sigma_0), \Lambda(\sigma_3)$ has no $\sigma_1,\sigma_2$ components.

    Use the condition that $\Lambda$ is single qubit incoherent channel, we conclude that it must has the form   \(\Lambda\left(\frac{1}{2}[\mathbb{I} + \mathbf{w} \cdot  \vec\sigma]\right) = \frac{1}{2}[\mathbb{I} + (\mathbf{t} + \Gamma\mathbf{w}) \cdot  \vec\sigma]\) with $\mathbf{t} = (0,0,0), \quad  \Gamma = \text{diag}(\Gamma_0,\Gamma_0,1)$.
\end{proof}
For noise channel $\Lambda_{\Gamma_0}$ on graph state that satisfied lemma~\ref{lemma:commute_channel}, It can directly operate on $\ket{+}$ state and output a mixed state
\begin{align}
    \Lambda_{\Gamma_0}(\ket{+}\!\bra{+}) = \frac{1 - \Gamma_0}{2} \ket{0}\!\bra{0} + \frac{1 - \Gamma_0}{2} \ket{1}\!\bra{1} + \Gamma_0 \ket{+}\!\bra{+}
\end{align}
This means it decoheres $\ket{+}$ state into states $\{\ket{0}\!\bra{0}, \ket{1}\!\bra{1}\}$  with an equal probability $p = {2^{-1}}{(1 - \Gamma_0)}$. These state will not entangled with the other single qubit states under the operation of following phase gates. When $\Gamma_0 = 0$, $\Lambda_{\Gamma_0}$ is the $Z$-projection channel $\Lambda_z(\rho) = ({1}/{2})\rho+({1}/{2})Z\rho Z$, which fully decohere the $\ket{+}$ state and equivalent to remove this vertex on the final graph state.
Based on this, many methods have been proposed to discuss the percolation behavior in the IQP-circuits~\cite{nelson2024polynomial,rajakumar2024polynomialtimeclassicalsimulationnoisy}. These methods find the component of $Z$-projection channel in Pauli noise. To be precise, we refer the theorem proved by Joel \emph{et.al.}~\cite{rajakumar2024polynomialtimeclassicalsimulationnoisy}. 
\begin{lemma} \label{lemma:z-projective_of_Paulinoise}
    Any Pauli channel $\mathcal{N}(\rho) = p \rho + p_x X \rho X + p_y Y \rho Y+ p_z Z \rho Z$ can be decomposed as
    \begin{align}
        \mathcal{N}(\rho) = (1 - 2p)\mathcal{N}^{(1)}(\rho) + 2p\mathcal{N}^{(2)} \circ \Lambda_z(\rho),
    \end{align}
    where $p = p_Z + \min(p_X,p_Y)$ , $N^{(2)}_i$ has parameter $(0,\frac{\min(p_x, p_y)}{p},0)$  , and  $N^{(1)}_i$ has parameter $(0,\frac{|p_x, p_y|}{1 - 2p},0)$ if $p_x < p_y$ and $(\frac{|p_x, p_y|}{1 - 2p},0,0)$ if $p_x \geq p_y$ .
\end{lemma}
Unfortunately, we show that similar decomposition cannot be done for amplitude damping channel.
\begin{lemma}\label{lemma: no_go_non-unital}
    It is impossible to decomposed the amplitude damping channel as
    \begin{align}
        \mathcal{N}(\rho) = (1 - p)\mathcal{N}^{(1)}(\rho) + p\mathcal{N}^{(3)} \circ \mathcal{N}^{(2)} \circ \Lambda_z(\rho),
    \end{align}
    where $\mathcal{N}^{(1)},\mathcal{N}^{(2)}$ is Pauli noise and  $\mathcal{N}^{(3)}$ is non-unital noise with z-shift only.
\end{lemma}
\begin{proof}
    If such a decomposition exist, we have
    \begin{align}
        \begin{pmatrix}
            1&0&0&0\\
            0&\sqrt{r}&0&0\\
            0&0&\sqrt{r}&0\\
            1-r&0&0&r
        \end{pmatrix} =
        p\begin{pmatrix}
            1&0&0&0\\
            0&\sqrt{r}/p&0&0\\
            0&0&\sqrt{r}/p&0\\
            0&0&0&\lambda
        \end{pmatrix}
        +(1-p)
        \begin{pmatrix}
            1&0&0&0\\
            0&0&0&0\\
            0&0&0&0\\
            q&0&0&m
        \end{pmatrix}
    \end{align}
    The second matrix has vanished entries on the second and the third column because it is a $Z$-projection channel followed by Pauli noise and  non-unital noise with z-shift only. Then we have
    \begin{align}
        1 - r = (1 - p)q, \quad r= p \lambda + (1- p)m 
    \end{align}
    To be a valid Pauli channel, $p > \sqrt{r}$ must be satisfied.
    However $q = {(1 - r)}{(1 - p)}^{-1} \leq 1$ implies that $ p < r$. This is conflict with the condition $p > \sqrt{r}$ in the case $0<r<1$.    
    \end{proof}
    Lemma~\ref{lemma: no_go_non-unital} means that our algorithm for amplitude sampling only applied for Pauli-noise. How to extend this algorithm to the case of non-unital work is still under consideration. 
    

\end{document}